%% preamble.tex — shared packages and macros
%% ─────────────────────────────────────────────────────────────────────────────


%% ── Core packages ────────────────────────────────────────────────────────────
\usepackage[utf8]{inputenc}        % UTF-8 source encoding
\usepackage{graphicx}              % figures (\includegraphics)
\usepackage{hyperref}              % hyperlinks in PDF
\usepackage{xspace}                % smart spacing after macros


%% ── Tables ───────────────────────────────────────────────────────────────────
\usepackage{booktabs}              % \toprule, \midrule, \bottomrule
\usepackage{multirow}              % multi-row table cells
\usepackage{diagbox}               % diagonal box in table headers
\usepackage[flushleft]{threeparttable} % table notes


%% ── Figures ──────────────────────────────────────────────────────────────────
\usepackage[caption=false]{subfig} % sub-figures


%% ── Code listings ────────────────────────────────────────────────────────────
\usepackage{listings}
\usepackage[htt]{hyphenat}         % hyphenation in \texttt

% YAML language definition for listings
\newcommand\YAMLcolonstyle{\color{red}\mdseries}
\newcommand\YAMLkeystyle{\color{black}\bfseries}
\newcommand\YAMLvaluestyle{\color{blue}\mdseries}

\makeatletter
\newcommand\language@yaml{yaml}
\expandafter\expandafter\expandafter\lstdefinelanguage
\expandafter{\language@yaml}
{
  keywords={true,false,null,y,n},
  keywordstyle=\color{darkgray}\bfseries,
  basicstyle=\YAMLkeystyle,
  sensitive=false,
  comment=[l]{\#},
  morecomment=[s]{/*}{*/},
  commentstyle=\color{purple}\ttfamily,
  stringstyle=\YAMLvaluestyle\ttfamily,
  moredelim=[l][\color{orange}]{\&},
  moredelim=[l][\color{magenta}]{*},
  morestring=[b]',
  morestring=[b]",
  literate = {---}{{\ProcessThreeDashes}}3
             {\ -\ }{{\mdseries\ -\ }}3,
}
\makeatother


%% ── Drawing / diagrams ───────────────────────────────────────────────────────
\usepackage{tikz}
\usepackage[framemethod=TikZ]{mdframed}


%% ── Colors ───────────────────────────────────────────────────────────────────
\usepackage{color, colortbl}
\usepackage{pifont}
\definecolor{lightgray}{gray}{0.9}


%% ── Misc formatting ──────────────────────────────────────────────────────────
\usepackage{pmboxdraw}             % box-drawing characters


%% ── To-do / review notes ─────────────────────────────────────────────────────
%% Swap the two \usepackage lines to hide all notes in the final version.
\usepackage[]{todonotes}           % show notes
% \usepackage[disable]{todonotes}  % hide notes

% Author notes — add/rename as needed
\newcommand{\authorA}[1]{\todo[inline,color=green!40]{\footnotesize{AUTHOR A: #1}}}
\newcommand{\authorB}[1]{\todo[inline,color=red!40]{\footnotesize{AUTHOR B: #1}}}
\newcommand{\authorC}[1]{\todo[inline,color=blue!40]{\footnotesize{AUTHOR C: #1}}}

% Reviewer notes (for revision rounds)
\newcommand{\reviewerOne}[1]{\todo[inline,color=olive!40]{\footnotesize{Reviewer 1: #1}}}
\newcommand{\reviewerTwo}[1]{\todo[inline,color=orange!40]{\footnotesize{Reviewer 2: #1}}}
\newcommand{\reviewerThree}[1]{\todo[inline,color=teal!40]{\footnotesize{Reviewer 3: #1}}}

% Review-response macros
\newcommand{\review}[1]{\par\noindent{\color{darkgray}{\textbf{Reviewer comment:} #1}}\smallskip}
\newcommand{\reply}[1]{\par\noindent{\color{blue}{\textbf{Authors' response:} #1}}\bigskip}
\newcommand{\reviewer}[2]{\todo[color=red!40, inline]{\footnotesize{\##1: #2}}}


%% ── Cross-reference shorthands ───────────────────────────────────────────────
\newcommand{\fig}[1]{Fig.~\ref{#1}}
\newcommand{\listing}[1]{Listing~\ref{#1}}
\newcommand{\tab}[1]{Table~\ref{#1}}
\newcommand{\sect}[1]{Section~\ref{#1}}


%% ── Latin / English abbreviations ────────────────────────────────────────────
\newcommand{\etal}{et al.\xspace}
\newcommand{\ie}{i.e.,\xspace}
\newcommand{\eg}{e.g.,\xspace}


%% ── Change-tracking macros ───────────────────────────────────────────────────
%% Use \changed{} to highlight revised text, or disable by making it a no-op.
%\newcommand{\changed}[1]{#1}
 \newcommand{\changed}[1]{{\color{red}#1}}


%% ── Custom callout box ───────────────────────────────────────────────────────
\global\mdfdefinestyle{takeawaybox}{%
  linewidth=1pt, linecolor=lightgray!80,
  backgroundcolor=lightgray!20,
  innerleftmargin=6pt, innerrightmargin=6pt,
  innertopmargin=4pt,  innerbottommargin=4pt,
  nobreak=true, roundcorner=5pt}
\newenvironment{takeawaybox}
  {\begin{mdframed}[style=takeawaybox]\emph{Takeaways.}}
  {\end{mdframed}\smallskip}


%% ── Inline quote helper ──────────────────────────────────────────────────────
\newcommand{\say}[1]{``\emph{#1}''}


%% ── Domain-specific shorthands (edit / remove as needed) ─────────────────────
% Tools / platforms
\newcommand{\github}{GitHub\xspace}
\newcommand{\gitlab}{GitLab\xspace}
\newcommand{\bitbucket}{BitBucket\xspace}
\newcommand{\git}{git\xspace}

% Languages / package managers
\newcommand{\python}{Python\xspace}
\newcommand{\js}{JavaScript\xspace}
\newcommand{\java}{Java\xspace}
\newcommand{\npm}{npm\xspace}

% Operating systems
\newcommand{\macos}{macOS\xspace}

% BibTeX logo
\def\BibTeX{{\rm B\kern-.05em{\sc i\kern-.025em b}\kern-.08em
    T\kern-.1667em\lower.7ex\hbox{E}\kern-.125emX}}
